\clearpage\fancyhead[L]{\small ANZEIGEN. LOGARITHMISCHE TAFELN.}
\SourceLeaf{262}
\pdfbookmark[1]{Babbage. 1828 Januar 19.}{babbage1828}\par\Needspace{10\baselineskip}\begin{center}\rule{0.7\linewidth}{0.4pt}\par Göttingische gelehrte Anzeigen. 1828 Januar 19.\par\rule{0.7\linewidth}{0.4pt}\end{center}

\noindent\textit{Table of logarithms of the natural numbers, from} 1 \textit{to} 108000, \textit{by} \textsc{Charles}
\textsc{Babbage.} \textit{Stereotyped. London. Printed for J. Mawman.} 1827. 202 S. gr. 8.

Dieser neue Druck der Logarithmentafeln zeichnet sich vor anderen durch eine bewusstere Beachtung kleiner Nebenumstände aus. Wer nur gelegentlich einige Logarithmen in den Tafeln nachschlagen muss, verlangt von ihnen hauptsächlich größtmögliche Richtigkeit. Für andere jedoch, denen die Tafeln ein tägliches Arbeitsgerät sind, werden selbst die geringfügigsten Umstände wichtig, die den bequemen Gebrauch beeinflussen können. Farbe, Stärke und Schönheit des Papiers; Format; Größe, Schärfe und gefällige Form der Drucktypen; Beschaffenheit der Druckerschwärze; Anordnung der Zahlen, damit man das Gesuchte schnell und sicher findet, ohne die Augen zu ermüden; Vorhandensein von allem, was man braucht, aber auch Abwesenheit von allem, was man nicht braucht und was die Übersicht nur stören würde: All diese Umstände erhalten eine gewisse Bedeutung bei einer Tätigkeit, die man täglich hundertmal wiederholt. Freilich hängt manches von der Eigenart und Gewohnheit des Rechnenden ab, sodass wohl nicht eine Ausgabe allen am besten gefallen kann. Für Kurzsichtige ist ein großes Format unbequem, und sie bevorzugen kleinere Drucktypen, während es bei Weitsichtigen umgekehrt ist. Weniger geübte Rechner schätzen diesen oder jenen Zusatz, den geübtere Rechner, die mehrere Teile der Operationen mühelos im Kopf ausführen, lieber weglassen würden, weil alles Überflüssige nur stört. Dennoch gibt es auch allgemein gültige Regeln. Der Herausgeber der vorliegenden Tafeln hat in der Vorrede ein Dutzend solcher Vorschriften zusammengestellt, die er aus dem Vergleich vieler Logarithmentafeln abgeleitet hat. Sie leuchten meist unmittelbar ein, obwohl einige wohl nur mit Einschränkungen anzuerkennen sein dürften.

Am meisten unterscheidet sich die gegenwärtige Ausgabe der Logarithmentafeln von anderen dadurch, dass sie auf farbiges (gelbes) Papier gedruckt ist. Man gewöhnt sich bald daran und findet es, besonders zum Gebrauch bei künstlichem Licht, \SourceLeaf{263} angenehmer als weißes. Die Anordnung ist im Wesentlichen die übliche. Nach einem Wechsel der dritten Ziffer ist bei den übrigen Logarithmen derselben Zeile die vierte Ziffer in anderen Drucktypen gesetzt, die etwa halb so groß wie die übrigen sind. Für den Referenten, der auf solche Warnzeichen überhaupt wenig Wert legt, haben diese Typen auch nach einigen Monaten des Gebrauchs noch nicht das fremdartig Störende verloren. Er würde die in \textsc{Vega}’s Tafel verwendeten Sternchen oder die Punkte vorziehen, die in die besseren Exemplare von \textsc{Taylor}’s Tafeln eingedruckt sind. Sternchen wollte der Herausgeber nicht verwenden, weil sonst die Spalten zu breit geworden wären und dann kein Raum geblieben wäre, um die als Sekunden aufgefassten Zahlen in Grade, Minuten und Sekunden umzuwandeln. Der Referent hätte diese Umrechnungsspalte, die genauso eingerichtet ist wie in \textsc{Callet}’s Tafeln, gern entbehrt. Ein geübter Rechner wird für eine so leicht im Kopf auszuführende Umrechnung nicht erst die Tafeln aufschlagen. Zweckmäßig ist aber, dass der Herausgeber diese Spalten wenigstens durch eine kräftige Linie von der Logarithmentafel getrennt hat. Schließlich ist der Fall, in dem die letzte Ziffer erhöht wurde, weil der fortgesetzte Logarithmus als achte Ziffer 5 oder mehr gehabt hätte, durch einen Punkt unter dieser siebten Ziffer gekennzeichnet. Diese Einrichtung ist wenigstens besser als die von \textsc{Prasse} gewählte Kennzeichnung durch anders geformte Drucktypen. Doch behalten auch so einige Ziffern noch etwas unangenehm Fremdartiges, das nach unserer schon bei der Anzeige der \textsc{Prasse}’schen Tafeln in diesen Blättern [1811 Mai 25] geäußerten Ansicht durch den Nutzen nicht aufgewogen wird.

Im Übrigen lassen Drucktypen und Papier bei dieser Ausgabe der Logarithmentafeln nichts zu wünschen übrig, und auf die Richtigkeit ist die größte Sorgfalt verwendet worden. Bei einer Ausgabe, die den kleinen, für den bequemen Gebrauch wichtigen Umständen so viel Sorgfalt widmet, fällt die spärliche Angabe der Proportionalteile auf: Häufig muss man zurück- oder vorblättern, um die gerade benötigten zu finden.
\SourceLeaf{264}
